\noindent\textbf{Dificultad:} Difícil (login)

\subsection{Descripción del nivel}
Pantalla de \emph{login}: se elige un usuario propio (mínimo 5 caracteres) y se debe ingresar la contraseña que el binario considera válida para ese usuario.

\subsection{Análisis estático}
% Cómo se ubicó la rutina de validación: cadenas, referencias cruzadas, funciones relevantes.
\textit{[A completar.]}

\subsection{Análisis dinámico}
% Breakpoints usados, registros/memoria observados en la comparación.
\textit{[A completar.]}

\subsection{Algoritmo de validación}
% Pseudocódigo del algoritmo reconstruido.
\textit{[A completar.]}

\subsection{Solución (keygen)}
\textit{[A completar.]}

\begin{lstlisting}[language=Python]
# Keygen - Nivel 3
# TODO: pegar el código del keygen
\end{lstlisting}

\subsection{Resultado}
Credencial generada y aceptada por el binario (ver Figura~\ref{fig:levels}):

\begin{lstlisting}[language={}]
Username: fierro
Password: C65715F2BE5918C4
[+] Level 3 PASSED
\end{lstlisting}
